\renewcommand{\Anno}{2003}

\renewcommand{\NumeroProblema}{3}

\renewcommand{\Testo}{\begin{enumerate}
	\item  Determinare la più piccola costante $a$ tale che
	$$6x^2+y^2+a \geq 4xy+y$$
	per ogni $x$ e $y$ reali.  Per tale valore di $a$ determinare le coppie $(x,y)$ di numeri reali per cui si ha uguaglianza.

	\item Determinare la più piccola costante $a$ tale che
	$$6x^2+y^2+a \geq 4xy+y$$
	per ogni $x$ e $y$ interi.  Per tale valore di $a$ determinare le coppie $(x,y)$ di interi per cui si ha uguaglianza.
\end{enumerate}}

\renewcommand{\Risposta}{\begin{enumerate}
		\item  $a=3/4$. Si ha uguaglianza se e solo se $(x,y)=(1/2,3/2)$. 
	
		\item  $a=0$. Si ha uguaglianza se e solo se $(x,y)$ è una delle seguenti 4 coppie: $(0,0)$, $(0,1)$, $(1,2)$, $(1,3)$.
\end{enumerate}}

\renewcommand{\Aiutino}{Portare la $a$ a destra e tutto il resto a
sinistra.
\begin{enumerate}
	\renewcommand{\labelenumi}{(\alph{enumi})} 
	\item Si tratta ora di determinare il minimo della funzione di 2 variabili scritta a sinistra: per fare questo si può cercare di scriverla come somma di quadrati (aggiungendo e togliendo quello che serve).

	\item Si tratta ora di determinare il minimo sugli interi della stessa funzione.  Tale minimo sarà intero e maggiore od uguale del minimo sui reali.  Una volta capito che il minimo è 0, occorre determinare tutti i punti a coordinate intere in cui la funzione vale 0, e per questo può tornare utile ancora una volta la scrittura della funzione come somma di quadrati.
\end{enumerate}}

\renewcommand{\SoluzioneUfficiale}{
}

\renewcommand{\SoluzioniAlternative}{
}

\renewcommand{\Osservazioni}{
}

\renewcommand{\Commenti}{
}
% Aggiunta redazionale: la soluzione preesistente non viene modificata.
\providecommand{\SoluzioneRedazionale}{}
\renewcommand{\SoluzioneRedazionale}{
% SSSUP-BOZZA-AI: SoluzioneRedazionale; autore=ChatGPT; data=2026-11-08; stato=DA-VERIFICARE
% SSSUP-PROVENIENZA: AI: elaborazione nuova dal testo
Portiamo a sinistra tutti i termini tranne la costante e poniamo
\[
f(x,y)=6x^2+y^2-4xy-y.
\]
La condizione richiesta è $f(x,y)+a\geq0$.

\paragraph{Punto (a).}
Il completamento dei quadrati dà
\[
f(x,y)=2\left(x-\frac12\right)^2+
\left(y-2x-\frac12\right)^2-\frac34.
\]
Il minimo reale è dunque $-3/4$, raggiunto solo in $(1/2,3/2)$. La minima costante è $a=3/4$ e l'uguaglianza si ha solo in quel punto.

\paragraph{Punto (b).}
Per gli interi è particolarmente utile la fattorizzazione
\[
f(x,y)=2x(x-1)+(y-2x)(y-2x-1).
\]
Ciascun prodotto di due interi consecutivi è non negativo. Quindi $f\geq0$ e il valore $a=0$ funziona; non si può scegliere $a<0$, perché $f(0,0)=0$. L'uguaglianza richiede $x\in\{0,1\}$ e $y-2x\in\{0,1\}$, cioè precisamente
\[
(0,0),\ (0,1),\ (1,2),\ (1,3).
\]
}
